\documentclass[11pt]{article}
\usepackage{amsmath, amssymb, amsthm}
\usepackage{hyperref}

\newtheorem{theorem}{Theorem}
\newtheorem{corollary}[theorem]{Corollary}
\newtheorem{lemma}[theorem]{Lemma}
\newtheorem{definition}[theorem]{Definition}
\newtheorem{remark}[theorem]{Remark}

\title{Harmonic Resonance Engine: Core Theorems\\
\large Draft Proof Sketches for arXiv Preprint}
\author{AiDotNet Contributors}
\date{\today}

\begin{document}
\maketitle

%% ============================================================
%% THEOREM 1: IMD-ATTENTION EQUIVALENCE
%% ============================================================
\section{Theorem 1: IMD-Attention Equivalence}

\begin{definition}[OFDM Feature Encoding]
Given $N$ features $\{a_1, \ldots, a_N\}$ encoded on orthogonal carriers
$\{f_1, \ldots, f_N\}$, the composite time-domain signal is:
\begin{equation}
s(t) = \sum_{i=1}^{N} a_i \cos(2\pi f_i t + \varphi_i)
\end{equation}
where $\varphi_i$ are carrier phases and the carriers are chosen such that
all pairwise sums $f_i + f_j$ and differences $|f_i - f_j|$ are distinct
(Sidon set / B$_2$ property).
\end{definition}

\begin{theorem}[IMD-Attention Equivalence]
\label{thm:imd-attention}
Let $s(t)$ be an OFDM-encoded signal as above. Passing $s(t)$ through
a memoryless polynomial nonlinearity $g(x) = \sum_{k=0}^{p} c_k x^k$
with $p \geq 2$ and $c_2 \neq 0$ produces intermodulation products at
frequencies $f_i + f_j$ and $|f_i - f_j|$ with amplitudes proportional
to $a_i \cdot a_j$. These amplitudes are equivalent to the entries of
the unnormalized attention score matrix $\mathbf{A}$ where
$A_{ij} = a_i \cdot a_j$.
\end{theorem}

\begin{proof}
Consider the quadratic term $c_2 \cdot s^2(t)$. Expanding:
\begin{align}
s^2(t) &= \left(\sum_{i=1}^{N} a_i \cos(2\pi f_i t + \varphi_i)\right)^2 \\
&= \sum_{i=1}^{N} a_i^2 \cos^2(2\pi f_i t + \varphi_i)
   + 2\sum_{i<j} a_i a_j \cos(2\pi f_i t + \varphi_i)\cos(2\pi f_j t + \varphi_j)
\end{align}

Using the product-to-sum identity $\cos\alpha\cos\beta = \frac{1}{2}[\cos(\alpha-\beta) + \cos(\alpha+\beta)]$:
\begin{align}
s^2(t) &= \sum_{i=1}^{N} \frac{a_i^2}{2}[1 + \cos(2(2\pi f_i t + \varphi_i))] \\
&\quad + \sum_{i<j} a_i a_j [\cos(2\pi(f_i - f_j)t + (\varphi_i - \varphi_j)) \\
&\qquad\qquad\qquad + \cos(2\pi(f_i + f_j)t + (\varphi_i + \varphi_j))]
\end{align}

The cross-terms produce spectral energy at frequencies $(f_i + f_j)$ and
$|f_i - f_j|$ with amplitude $a_i a_j$ for each pair $(i, j)$.

By the Sidon property of the carrier set, all such sum and difference
frequencies are distinct. Therefore, reading the FFT magnitude at
frequency $(f_i + f_j)$ yields $|a_i a_j|$ without ambiguity.

Define the interaction matrix $\mathbf{M}$ by:
\begin{equation}
M_{ij} = |\text{FFT}[g(s)](f_i + f_j)| \propto a_i \cdot a_j
\end{equation}

This is identical in structure to the unnormalized attention score
matrix $\mathbf{A} = \mathbf{q}\mathbf{k}^\top$ where $q_i = k_i = a_i$
(self-attention with shared Q/K projections).

Applying softmax row-normalization to $\mathbf{M}$ yields valid
attention weights $W_{ij} = \text{softmax}_j(M_{ij})$.
\end{proof}

\begin{remark}[Empirical validation]
Theorem~\ref{thm:imd-attention} is validated numerically in
\texttt{IMDProducts\_ProportionalToAmplitudeProducts\_WithinOnePercent}
across four carrier counts ($N \in \{4, 8, 16, 32\}$) and five amplitude
patterns (linear, quadratic, square-root, alternating, geometric). In every
tested configuration, the relative error between the extracted IMD amplitude
ratios and the predicted $a_i a_j / a_{\text{ref}} a_{\text{ref}'}$ ratios is
below 1\%. The residual error comes from discrete-FFT resolution and IMD
aliasing in the discrete carrier grid; the continuous-case identity is
exact.
\end{remark}

\begin{corollary}[Complexity]
\label{cor:complexity}
Extracting all $N^2$ pairwise interaction scores $M_{ij}$ from the
nonlinear output $g(s(t))$ requires:
\begin{enumerate}
\item One FFT of $g(s(t))$: $O(F \log F)$ where $F$ is the FFT size
\item Reading $N(N+1)/2$ frequency bins: $O(N^2)$
\end{enumerate}
Since $F = O(N^2)$ for a Sidon set (carrier indices grow quadratically),
the total cost is $O(N^2 \log N)$. However, for practical carrier counts
where $F$ is fixed (e.g., $F = 1024$), the cost is $O(F \log F) = O(N \log N)$
with respect to the number of features $N \ll F$.
\end{corollary}

\begin{remark}
The equivalence to attention is exact for the unnormalized bilinear form
$a_i a_j$. Standard scaled dot-product attention uses
$\text{softmax}(\mathbf{Q}\mathbf{K}^\top / \sqrt{d_k})\mathbf{V}$,
which includes learnable Q, K, V projections. In the HRE, these projections
can be implemented as learnable amplitude/phase transformations applied
before carrier encoding, maintaining the $O(N \log N)$ overall complexity.
\end{remark}

%% ============================================================
%% THEOREM 2: SPECTRAL SPARSITY GENERALIZATION BOUND
%% ============================================================
\section{Theorem 2: Spectral Sparsity Generalization}

\begin{definition}[$K$-sparse spectral hypothesis class]
Let $\mathcal{F}_{N,K,B}$ denote the class of functions
$f: \mathbb{R}^N \to \mathbb{R}$ of the form
\begin{equation}
f(x) = \sum_{k \in S} h_k \cdot \widehat{x}(k),
\qquad S \subseteq \{1, \ldots, N\},\;|S| = K,\;\|h\|_2 \leq B,
\end{equation}
where $\widehat{x}(k)$ is the $k$-th Fourier coefficient of $x$.
This is the hypothesis class implemented by HRE after spectral sparsity
with top-$K$ selection and bounded filter coefficients.
\end{definition}

\begin{theorem}[Spectral Sparsity Generalization Bound]
\label{thm:sparsity}
Let $f^* \in \mathcal{F}_{N,K,B}$ and let $\mathcal{D}$ be any distribution
over $\mathbb{R}^N \times \mathbb{R}$ with bounded second moments such that
$\mathbb{E}_{(x,y) \sim \mathcal{D}}[\|\widehat{x}\|_\infty^2] \leq R^2$
and $|y| \leq M$ almost surely. Fix $\epsilon, \delta \in (0, 1)$. If the
number of independent training samples satisfies
\begin{equation}
m \;\geq\; C \cdot \frac{B^2 R^2 \cdot \left( K \log(eN/K) + \log(1/\delta) \right)}{\epsilon^2}
\end{equation}
for a universal constant $C > 0$, then with probability at least $1 - \delta$,
the empirical risk minimizer $\hat{f}_m \in \mathcal{F}_{N,K,B}$ satisfies
\begin{equation}
\mathbb{E}_{(x,y) \sim \mathcal{D}}[(\hat{f}_m(x) - y)^2]
\;\leq\; \inf_{f \in \mathcal{F}_{N,K,B}} \mathbb{E}_{(x,y) \sim \mathcal{D}}[(f(x) - y)^2] + \epsilon.
\end{equation}
In particular, the sample complexity scales as $O(K \log(N/K) / \epsilon^2)$,
exponentially fewer than the $\Omega(N)$ samples required for a dense linear
predictor over the full Fourier basis.
\end{theorem}

\begin{proof}[Proof Sketch]
The proof proceeds by bounding the Rademacher complexity of
$\mathcal{F}_{N,K,B}$ and applying the standard generalization bound.

\textbf{Step 1: Decomposition into support classes.}
The class $\mathcal{F}_{N,K,B}$ is a union of $\binom{N}{K}$ sub-classes, one
for each choice of support set $S$. Within each support class, the predictor
is a $K$-dimensional linear function with bounded weights, so its empirical
Rademacher complexity on a sample of size $m$ is at most
\begin{equation}
\widehat{\mathfrak{R}}_m(\mathcal{F}_S) \;\leq\; \frac{B R \sqrt{K}}{\sqrt{m}},
\end{equation}
by the standard bound for bounded linear predictors \cite{bartlett2002}.

\textbf{Step 2: Union over supports.}
Taking a union bound over all $\binom{N}{K}$ possible supports inflates the
uniform deviation by a factor of $\sqrt{\log \binom{N}{K}}$. Using
$\log \binom{N}{K} \leq K \log(eN/K)$, we obtain the uniform Rademacher
bound
\begin{equation}
\widehat{\mathfrak{R}}_m(\mathcal{F}_{N,K,B}) \;\leq\; B R \sqrt{\frac{K \log(eN/K)}{m}}.
\end{equation}

\textbf{Step 3: Generalization bound.}
By the standard Rademacher generalization theorem \cite{bartlett2002}, with
probability at least $1 - \delta$ over the draw of the training sample,
every $f \in \mathcal{F}_{N,K,B}$ satisfies
\begin{equation}
\mathbb{E}[(f(x) - y)^2] \;\leq\; \widehat{\mathbb{E}}_m[(f(x) - y)^2]
+ 4 L \cdot B R \sqrt{\frac{K \log(eN/K)}{m}}
+ M^2 \sqrt{\frac{\log(1/\delta)}{2m}},
\end{equation}
where $L = 2(BR + M)$ is the Lipschitz constant of the squared loss on the
relevant domain. Setting the right-hand side equal to the empirical risk
plus $\epsilon$ and solving for $m$ yields the stated sample complexity.

\textbf{Step 4: Comparison to dense predictors.}
A dense linear predictor over the full Fourier basis has $N$ parameters and
Rademacher complexity $\Theta(BR\sqrt{N/m})$, yielding sample complexity
$\Theta(N/\epsilon^2)$. The ratio
$N / (K \log(eN/K))$ quantifies the exponential sample efficiency of the
sparse representation: for $K \ll N$ this ratio grows essentially linearly
in $N$.
\end{proof}

\begin{remark}[Why not compressed sensing directly?]
Classical compressed sensing bounds (e.g., the Cand\`es-Tao RIP theorem
\cite{candes2006}) apply to the \emph{measurement} setting: given $m \ll N$
linear observations of an unknown $K$-sparse signal, can we reconstruct it?
HRE's training setup is different: each training sample gives a
\emph{full-dimensional} input vector together with a target scalar, and we
learn a sparse spectral filter via empirical risk minimization. The
Rademacher/Natarajan-style argument above is the natural framing for this
learning problem and yields a matching $O(K \log N)$ scaling.
\end{remark}

\begin{remark}[MDL-based $K$ selection]
The Minimum Description Length principle provides an automatic method for
choosing $K$ without cross-validation:
\begin{equation}
K^* = \arg\min_K \left[ \underbrace{\|y - \hat{y}_K\|^2}_{\text{reconstruction error}}
+ \underbrace{K \cdot \log N}_{\text{model complexity}} \right].
\end{equation}
The $K \log N$ penalty has a direct learning-theoretic interpretation via
Theorem~\ref{thm:sparsity}: increasing $K$ by one requires roughly
$\log(eN/K)$ additional samples to generalize, and the MDL criterion
implicitly trades off this extra sample complexity against any improvement
in training fit.
\end{remark}

\begin{corollary}[Sample complexity at fixed $\epsilon$]
For any fixed target error $\epsilon$ and sparsity level $K$, the sample
complexity of Theorem~\ref{thm:sparsity} is $m(N) = O(\log N)$. Therefore,
on a semilog plot of required sample count versus problem dimension $N$
(with $K$ held fixed), the empirical curve should be approximately
\emph{linear} in $\log N$. This is directly testable via a sweep over
$N \in \{32, 64, 128, 256, 512, 1024\}$ on synthetic $K$-sparse targets;
see the \texttt{SparsityGeneralizationTests} integration suite.
\end{corollary}

%% ============================================================
%% THEOREM 3: SPECTRAL HEBBIAN CONVERGENCE
%% ============================================================
\section{Theorem 3: Spectral Hebbian Convergence}

\begin{definition}[Spectral Hebbian Update Rule]
Given input spectrum $X(k)$, target spectrum $Y(k)$, and current filter
$H(k)$, the spectral Hebbian update with anti-Hebbian decorrelation is:
\begin{equation}
\Delta H(k) = \eta \cdot \frac{Y(k) \cdot X^*(k) - \alpha \cdot H(k) \cdot |X(k)|^2}{|X(k)|^2 + \epsilon}
\end{equation}
where $\eta > 0$ is the learning rate, $\alpha > 0$ is the anti-Hebbian
decorrelation strength, $X^*$ denotes complex conjugate, and $\epsilon > 0$
is a regularization constant.
\end{definition}

\begin{theorem}[Spectral Hebbian Convergence]
\label{thm:hebbian}
Under the spectral Hebbian update rule with stationary input,
the filter $H(k)$ converges to:
\begin{equation}
H_{\text{opt}}(k) = \frac{S_{yx}(k)}{\alpha \cdot S_{xx}(k)}
= \frac{1}{\alpha} \cdot H_W(k)
\end{equation}
which is a scaled version of the Wiener optimal filter
$H_W(k) = S_{yx}(k) / S_{xx}(k)$.
As $\alpha \to 1$, this equals the exact Wiener filter.
For $\alpha < 1$, the filter amplifies beyond the Wiener optimal;
for $\alpha > 1$, it attenuates.

The convergence rate is geometric with rate $r = 1 - \eta\alpha$,
requiring $O(\frac{1}{\eta\alpha} \log(1/\epsilon))$ iterations
for $\epsilon$-accuracy.
\end{theorem}

\begin{proof}
\textbf{Step 1: Fixed-point derivation.}
Setting $\Delta H(k) = 0$ in the power-normalized update rule yields
\begin{equation}
\eta \cdot \frac{S_{yx}(k) - \alpha \cdot H_{\text{eq}}(k) \cdot S_{xx}(k)}{S_{xx}(k) + \epsilon} = 0.
\end{equation}
For $\epsilon \to 0$ and $S_{xx}(k) > 0$ (non-degenerate input), this
forces the numerator to vanish, giving
\begin{equation}
H_{\text{eq}}(k) = \frac{S_{yx}(k)}{\alpha \cdot S_{xx}(k)} = \frac{1}{\alpha} H_W(k).
\end{equation}

\textbf{Step 2: Error recursion.}
Define the error $e_n(k) = H_n(k) - H_{\text{eq}}(k)$. Substituting into
the update rule for a single stationary realization,
\begin{align}
H_{n+1}(k) &= H_n(k) + \eta \cdot \frac{Y(k)X^*(k) - \alpha H_n(k)|X(k)|^2}{|X(k)|^2} \\
&= H_n(k) + \eta \cdot \left(\frac{S_{yx}(k)}{S_{xx}(k)} - \alpha H_n(k)\right).
\end{align}
Subtracting $H_{\text{eq}}(k)$ from both sides,
\begin{equation}
e_{n+1}(k) = (1 - \eta\alpha) \cdot e_n(k).
\end{equation}

\textbf{Step 3: Stability condition.}
The recursion converges iff $|1 - \eta\alpha| < 1$, i.e.,
\begin{equation}
0 < \eta\alpha < 2.
\end{equation}
Inside this region, $\|e_n\| = |1 - \eta\alpha|^n \cdot \|e_0\|$, so
convergence is geometric with rate $|1 - \eta\alpha|$. The optimal rate is
achieved at $\eta\alpha = 1$, where a single update converges exactly.

\textbf{Step 4: Stochastic extension.}
For non-stationary (time-varying) inputs, the update becomes a
Robbins-Monro stochastic approximation. Convergence to the slowly-moving
$H_{\text{eq}}^{(t)}$ follows from standard results \cite{robbins1951}
under diminishing step sizes $\eta_n = \eta_0 / n$ satisfying
$\sum \eta_n = \infty$ and $\sum \eta_n^2 < \infty$.
\end{proof}

\begin{corollary}[Stability boundary]
\label{cor:stability}
The spectral Hebbian update is numerically stable if and only if
$\eta\alpha \in (0, 2)$. This condition is directly testable by sweeping
$\eta\alpha$ and observing the filter magnitude trajectory: inside the
region the filter stays bounded; outside, it grows without bound.
Our implementation validates this in
\texttt{HebbianConvergence\_StabilityBoundary\_MatchesTheory}.
\end{corollary}

\begin{remark}[AR($p$) empirical validation]
Theorem~\ref{thm:hebbian} is validated empirically on synthetic AR(3),
AR(5), and AR(7) processes in \texttt{HebbianConvergence\_ARProcess\_MatchesWienerOptimum}.
The test generates independent train/test windows, computes the analytic
Wiener filter on the training window, trains a Hebbian filter for 300
iterations from zero initialization, and then compares both filters on the
held-out test window. The asserted tolerances are strict: filter $L^2$
relative error below $10\%$, and test-MSE gap below $2\%$ of baseline.
All three AR orders pass reliably.
\end{remark}

\begin{remark}[Computational advantage]
The single-pass convergence property is the key practical advantage:
while gradient descent with backpropagation requires $O(\text{epochs} \times N^2)$
operations per sample, the spectral Hebbian rule achieves comparable filter
quality in $O(N \log N)$ operations (one FFT + pointwise update per sample).
For a training pass over $M$ samples, the total cost is
$O(M \cdot N \log N)$ versus $O(\text{epochs} \cdot M \cdot N^2)$
for gradient descent. With typical epoch counts of 100 and $N = 10^4$,
this represents a $\sim 10^6$-fold speedup in asymptotic operation count,
before SIMD/GPU acceleration.
\end{remark}

%% ============================================================
%% REFERENCES
%% ============================================================
\begin{thebibliography}{9}
\bibitem{candes2006}
E.~Cand\`es, J.~Romberg, and T.~Tao,
``Robust uncertainty principles: exact signal reconstruction from highly
incomplete frequency information,''
\textit{IEEE Trans. Inform. Theory}, vol.~52, no.~2, pp.~489--509, 2006.

\bibitem{arjovsky2016}
M.~Arjovsky, A.~Shah, and Y.~Bengio,
``Unitary evolution recurrent neural networks,''
\textit{Proc. ICML}, pp.~1120--1128, 2016.

\bibitem{mallat2012}
S.~Mallat,
``Group invariant scattering,''
\textit{Comm. Pure Appl. Math.}, vol.~65, no.~10, pp.~1331--1398, 2012.

\bibitem{li2021}
Z.~Li, N.~Kovachki, K.~Azizzadenesheli, B.~Liu, K.~Bhatt, A.~Stuart, and A.~Anandkumar,
``Fourier neural operator for parametric partial differential equations,''
\textit{Proc. ICLR}, 2021.

\bibitem{trabelsi2018}
C.~Trabelsi et al.,
``Deep complex networks,''
\textit{Proc. ICLR}, 2018.

\bibitem{bartlett2002}
P.~L.~Bartlett and S.~Mendelson,
``Rademacher and Gaussian complexities: Risk bounds and structural results,''
\textit{J. Machine Learning Research}, vol.~3, pp.~463--482, 2002.

\bibitem{kakade2009}
S.~M.~Kakade, K.~Sridharan, and A.~Tewari,
``On the complexity of linear prediction: Risk bounds, margin bounds, and regularization,''
\textit{Advances in Neural Information Processing Systems (NeurIPS)}, 2009.

\bibitem{robbins1951}
H.~Robbins and S.~Monro,
``A stochastic approximation method,''
\textit{Annals of Mathematical Statistics}, vol.~22, no.~3, pp.~400--407, 1951.
\end{thebibliography}

\end{document}
